\begin{figure*}[!hbt]
\centering

\begin{tikzpicture}[
    scale=0.70,
    transform shape,
    >=Latex,
    font=\normalsize,
    stage/.style={
        draw,
        rounded corners=5pt,
        minimum width=2.55cm,
        minimum height=1.0cm,
        align=center,
        inner sep=2pt,
        font=\small
    },
    boundary/.style={
        circle,
        draw,
        thick,
        minimum size=7mm,
        inner sep=0.5pt,
        font=\small\bfseries
    },
    btitle/.style={
        align=center,
        text width=2.65cm,
        font=\small\bfseries
    },
    bdetail/.style={
        align=center,
        text width=2.70cm,
        font=\scriptsize
    },
    side/.style={
        align=right,
        font=\small\bfseries,
        text width=1.9cm
    },
    mapcell/.style={
        draw=gray!55,
        align=center,
        minimum height=7mm,
        inner sep=2pt,
        font=\scriptsize
    }
]

% =================================================
% Colors
% =================================================
\definecolor{cpu}{RGB}{199,199,245}
\definecolor{qpu}{RGB}{224,194,220}

\definecolor{bOne}{RGB}{144,183,190}
\definecolor{bTwo}{RGB}{250,213,165}
\definecolor{bThree}{RGB}{247,231,164}
\definecolor{bFive}{RGB}{190,224,192}
\definecolor{bFour}{RGB}{184,220,232}
\definecolor{bSix}{RGB}{111,193,225}

% =================================================
% Generic HQC workflow
% =================================================

\node[stage,fill=cpu] (s1) at (0,0)
    {problem/data\\preparation};

\node[stage,fill=cpu] (s2) at (3,0)
    {execution\\planning};

\node[stage,fill=cpu] (s3) at (6,0)
    {compile\\+ bind};

\node[stage,fill=qpu] (s4) at (9,0)
    {external QPU\\work};

\node[stage,fill=cpu] (s5) at (12,0)
    {aggregate\\+ update};

\node[stage,fill=cpu] (s6) at (15,0)
    {program evolution\\(optional)};

\draw[->,thick] (s1) -- (s2);
\draw[->,thick] (s2) -- (s3);
\draw[->,thick] (s3) -- (s4);
\draw[->,thick] (s4) -- (s5);
\draw[->,thick] (s5) -- (s6);

% =================================================
% Legend
% =================================================

\begin{scope}[xshift=-2.2cm,yshift=-0.4cm]

\draw[rounded corners=2pt]
    (-0.20,1.70) rectangle (7.60,2.40);

\draw[fill=cpu]
    (0.05,1.92) rectangle (0.42,2.18);

\node[anchor=west,font=\small]
    at (0.55,2.05)
    {classical processing};

\draw[fill=qpu]
    (3.45,1.92) rectangle (3.82,2.18);

\node[anchor=west,font=\small]
    at (3.85,2.05)
    {quantum (QPU) processing};

\end{scope}

% =================================================
% Hybrid classical--quantum loop
% =================================================

\draw[->,thick,gray]
    (12,0.62)
    .. controls (11.2,1.35) and (4.0,1.35) ..
    (3,0.62);

\node[
    font=\small,
    text=gray
] at (7.6,1.55)
    {hybrid classical--quantum loop};

% =================================================
% Left-side labels
% =================================================

\node[side] at (-2.4,0)
    {HQC\\workflow};

\node[side] at (-2.4,-1.75)
    {Semantic\\restart\\boundaries};

% =================================================
% Semantic restart boundaries
%
% IMPORTANT:
% IDs denote invariant types, NOT chronological stage numbers.
% Thus B5 naturally occurs before the B4 iteration commit.
% =================================================

\node[boundary,fill=bOne] (b1)
    at (0,-1.65) {$B_1$};

\node[boundary,fill=bTwo] (b2)
    at (3,-1.65) {$B_2$};

\node[boundary,fill=bThree] (b3)
    at (6,-1.65) {$B_3$};

\node[boundary,fill=bFive] (b5)
    at (9,-1.65) {$B_5$};

\node[boundary,fill=bFour] (b4)
    at (12,-1.65) {$B_4$};

\node[boundary,fill=bSix] (b6)
    at (15,-1.65) {$B_6$};

\draw[dashed,thick] (s1.south) -- (b1.north);
\draw[dashed,thick] (s2.south) -- (b2.north);
\draw[dashed,thick] (s3.south) -- (b3.north);
\draw[dashed,thick] (s4.south) -- (b5.north);
\draw[dashed,thick] (s5.south) -- (b4.north);
\draw[dashed,thick] (s6.south) -- (b6.north);

% =================================================
% Boundary meaning:
%   1. work protected
%   2. evidence gained
% =================================================

% -----------------------
% B1
% -----------------------

\node[btitle]
    at (0,-2.40)
    {Problem\\State};

\node[bdetail]
    at (0,-3.30)
    {\textbf{Protects:} problem/ data preparation\\
     \textbf{Evidence:\\} semantic identity};

% -----------------------
% B2
% -----------------------

\node[btitle]
    at (3,-2.40)
    {Execution\\Plan};

\node[bdetail]
    at (3,-3.30)
    {\textbf{Protects:} grouping, batching, planning\\
     \textbf{Evidence:\\} execution/shot plan};

% -----------------------
% B3
% -----------------------

\node[btitle]
    at (6,-2.40)
    {Device\\Binding};

\node[bdetail]
    at (6,-3.30)
    {\textbf{Protects:\\} compilation and mapping\\
     \textbf{Evidence:\\} portability};

% -----------------------
% B5
% -----------------------

\node[btitle]
    at (9,-2.40)
    {Partial\\External Work};

\node[bdetail]
    at (9,-3.30)
    {\textbf{Protects:} completed\\QPU work\\
     \textbf{Evidence:} progress ledger / exact reuse};

% -----------------------
% B4
% -----------------------

\node[btitle]
    at (12,-2.40)
    {Iteration\\Commit};

\node[bdetail]
    at (12,-3.30)
    {\textbf{Protects:} completed optimizer/training update\\
     \textbf{Evidence:\\} continuation state};

% -----------------------
% B6
% -----------------------

\node[btitle]
    at (15,-2.40)
    {Program\\Evolution};

\node[bdetail]
    at (15,-3.30)
    {\textbf{Protects:} adaptive program construction\\
     \textbf{Evidence:} program identity/history};

% =================================================
% Important semantic note
% =================================================

\node[
    font=\small\itshape,
    text=gray!70!black,
    align=center
] at (7.5,-4.10)
{
Boundary IDs denote reusable invariant types, not chronological stage
numbers; therefore $B_5$ can occur before the $B_4$ iteration commit.
};

% =================================================
% Decision-role callout
% =================================================

\node[
    draw=gray!50,
    rounded corners=3pt,
    fill=gray!7,
    font=\small,
    align=center,
    text width=17.2cm,
    inner sep=3pt
] at (7.5,-4.80)
{
\textbf{Decision roles:}
$B_1$ gates semantic identity;
$B_3$ enables portability checks;
$B_5$ enables exact reuse/reissue accounting;
$B_4$ supports stable iterative re-entry; and
$B_6$ preserves evolving program identity.
};
\end{tikzpicture}

{\scriptsize
\setlength{\tabcolsep}{3pt}
\renewcommand{\arraystretch}{0.95}

\begin{tabularx}{\textwidth}{
    >{\raggedright\arraybackslash\bfseries}p{1.50cm}
    >{\centering\arraybackslash}X
    >{\centering\arraybackslash}X
    >{\centering\arraybackslash}X
    >{\centering\arraybackslash}X
    >{\centering\arraybackslash}X
    >{\centering\arraybackslash}X
}
\toprule
Workload
&
\shortstack{$B_1$ (Problem)}
&
\shortstack{$B_2$ (Plan)}
&
\shortstack{$B_3$ (Binding)}
&
\shortstack{$B_4$ (Iteration)}
&
\shortstack{$B_5$ (Partial work)}
&
\shortstack{$B_6$ (Evolution)}
\\
\midrule

VQE
&
Hamiltonian
&
groups/shots
&
compiled ansatz
&
optimizer state
&
measured groups
&
-- 
\\

ADAPT-VQE
&
Hamiltonian
&
groups/shots
&
compiled ansatz
&
optimizer state
&
measured groups
&
operator history
\\

QAOA
&
graph/cost
&
eval./shot plan
&
compiled QAOA
&
optimizer state
&
completed evals.
&
--
\\

VQC
&
dataset/encoding
&
batch/shot plan
&
compiled classifier
&
training state
&
samples/batches
&
--
\\

\bottomrule
\end{tabularx}
}

\caption{
\textbf{Semantic restart-boundary taxonomy and workload mapping.}
Semantic restart boundaries are defined by reusable workflow invariants
rather than algorithm-specific stage names. Each boundary protects a
different prefix of work and exposes evidence used in subsequent
restart decisions. VQE, ADAPT-VQE, QAOA, and the targeted VQC
instantiate the same invariant types differently; workloads may omit
or repeat boundary types. Boundary identifiers denote semantic
invariant types rather than chronological stage numbers, so the
partial-external-work boundary $B_5$ can occur before the
iteration-commit boundary $B_4$.
}
\label{fig:semantic_boundary_taxonomy}
\end{figure*}